% macros generadas; no editar a mano
\newcommand{\Nrep}{2}
\newcommand{\HrmsMean}{0.4699}
\newcommand{\HrmsStd}{0.0747}
\newcommand{\HeqMean}{0.4706}
\newcommand{\EoverN}{-7.939}
\newcommand{\AccMove}{39.2}
\newcommand{\AccArea}{12.3}
\newcommand{\LxMean}{149.52}
\newcommand{\LyMean}{155.38}
\newcommand{\Szero}{2.689}
\newcommand{\NsampTot}{598}
\newcommand{\RdfPeak}{1.445}
\newcommand{\Qmin}{0.040}
\newcommand{\Qripple}{0.090}
\newcommand{\GqBaja}{0.40}
\newcommand{\GqMedia}{0.75}
\newcommand{\MaximoFrase}{En la ventana $q\in[0.06,0.25]$~\AA$^{-1}$ el promedio de $G(q)$ tiene un m\'aximo ancho cerca de $q\approx 0.06$--$0.09$~\AA$^{-1}$, el lugar donde el paper lo espera para muestras m\'as largas que $70$~\AA. La dispersi\'on entre las dos r\'eplicas es grande a ese $q$, as\'i que no se identifica con el pico agudo que el paper atribuye al canal de conjugaci\'on de LCBOPII.}
\newcommand{\HrmsFrase}{El paper describe, para $N=8640$, una altura t\'ipica $\sim 0.6$~\AA\ (Fig.~1). Aqu\'i $h_{\mathrm{rms}}=0.4699$~\AA, que es $\sqrt{\langle h^2\rangle}$ y no el mismo estimador visual. Queda en el mismo orden de magnitud.}
